\section{Discussion}

\subsection{Similarity ranking and model-level interpretation}

Cell-line-to-tumour similarity ranking achieved top-1 same-lineage matching
accuracy of 54/56 biological cell-line groups (96.4\%; 95\% Wilson CI,
87.9--99.0\%) (Fig.~\ref{fig:cellline_centred_tumour_class_similarity}).
This result was computed after grouping replicate libraries into 56 biological
cell-line groups; the complementary patient-to-cell-line analysis ranked 58
candidate cell-line models for each patient tumour. The cell-line-to-tumour
result supports lineage-level transcriptomic fidelity, whereas the
patient-to-cell-line ranking provides the practical model-prioritisation view.

This bidirectional design addresses two complementary ranking questions.
The patient-to-cell-line direction ranks candidate cell-line models for each
patient tumour, yielding rank distributions across the cohorts. The
cell-line-to-tumour direction ranks patient tumour contexts for each biological
cell-line group and yielded the top-1 same-lineage matching accuracy reported
above. Patient-to-cell-line ranking should therefore be interpreted as model
prioritisation rather than functional validation. The broader breast cancer
ranking pattern is consistent with greater cohort size and within-lineage
heterogeneity. Same-lineage rank concentration also differed by cohort, with
retinoblastoma showing tighter similarity ranking than breast cancer. This cohort
dependence means that top-ranked lineage agreement should not be interpreted as
equivalent model-selection confidence across diseases. The confidence-score analysis further shows that
correct top-ranked lineage assignment does not imply equal match strength
across cohorts; downstream model selection should therefore consider both rank
position and the supporting similarity score
(Fig.~\ref{fig:cellline_centred_confidence_scores}).

Beyond the aggregate matching result, the framework identified model-specific
patterns that were consistent across ranking and graph analyses. The two
cell-line-to-tumour top-1 exceptions were also graph-resolved atypical
neuroblastoma profiles (Fig.~\ref{fig:per_cellline_feature_distance_combined}b).
Taken together, these recurrent atypical placements indicate that model
selection should consider transcriptomic neighbourhood structure in addition to
nominal lineage labels.

The graph framework also identified positive model-selection signals.
Bridge-like or high-betweenness anchors occupied central positions within
patient-referenced neighbourhood structure and may serve as candidate reference
models for component-level interpretation. In contrast, graph isolates provided
context-specific cautionary cases rather than lineage-label-based selection
alone; the full component annotation table reports the remaining anchors and
isolates in \suppTableRef{tab:supp_component_annotations}. The multiple resolved
breast cancer components are consistent with the well-established intrinsic
molecular heterogeneity of breast cancer
\citep{Parker2009PAM50,Perou2000Portraits} and could motivate PAM50-aware or
receptor-status-stratified validation rather than a single lineage-level
interpretation.


The 167-profile pan-cancer cell-line network extends this model-level view. Its
single connected component and high lineage assortativity suggest that the
cell-line profiles occupied structured but connected transcriptomic regions
rather than fully disconnected lineage-specific spaces, with HEME profiles
functioning as an auxiliary comparator lineage rather than a primary validation
target. Beyond individual-model atypical placement, the NBL--RBL community
co-assignment indicates incomplete lineage separability in the curated feature
space, and the assignment of GI\_ME\_N and CCLF\_PEDS\_0051\_T to a
BRCA-majority community shows that cross-lineage discordance was not restricted
to retinoblastoma. The separation of haematological comparator profiles into
multiple high-purity communities is consistent with substantial transcriptomic
subdivision within the auxiliary HEME set and reinforces that these profiles
should be interpreted as comparators rather than primary validation lineages.
This NBL--RBL graph adjacency is consistent with overlapping embryonal neural,
developmental, and proliferative transcriptomic programmes rather than a shared
terminal lineage, reflecting developmental states described in neuroblastoma
sympathoblasts and retinoblastoma cone photoreceptor precursors
\citep{Jansky2021NeuroblastomaDevelopment,Yang2021RetinoblastomaSingleCell}
(Fig.~\ref{fig:network}). This interpretation is limited to transcriptomic
co-assignment, because the analysis does not distinguish shared lineage biology,
cell-line-specific adaptation, profile or library effects, or limitations of
bulk transcriptomic similarity.

\subsection{Feature--distance sensitivity and molecular context}

These biological patterns depended on feature--distance choices and graph-based
resolution. Feature-selection performance varied with tumour lineage:
dispersion-based features under the Euclidean metric attained the highest
consensus fractions in breast cancer and retinoblastoma, whereas
co-expression-network-derived features under correlation distance attained the
highest consensus fractions in neuroblastoma. This representation dependence
indicates that no single feature--distance representation generalised uniformly
across lineages (Fig.~\ref{fig:feature_distance_combined}). Representation
choice was therefore consequential rather than a minor tuning parameter: the
best and weakest representations differed substantially within cohorts, and the
preferred metric pattern inverted across lineages. The
Euclidean-versus-correlation inversion across lineages suggests that the
dominant similarity signal differed by cohort: absolute expression differences
were more informative in breast cancer and retinoblastoma, whereas relative
co-expression structure was more informative in neuroblastoma.
Graph-based resolution organised the patient-referenced similarity graph into
connected components, isolates, and high-betweenness anchor nodes
(Fig.~\ref{fig:nbl_resolved_components}).

The molecular context for the graph and ranking patterns was obtained through
differential expression analysis, recurrence filtering, isolate-associated
marker rescue, and enrichment analysis. Recurrence filtering retained genes
that recurred above threshold across independent differential expression
contrasts, while isolate-associated marker rescue separately retained
additional genes associated with isolated profiles. Recurrence filtering
produced a smaller recurrent core, especially in retinoblastoma, where
only one recurrent-core gene remained. Isolate-associated marker rescue helped
preserve model-specific signal that strict recurrence criteria would otherwise
discard; this matters because smaller lineages and isolated graph structures
provide fewer independent contrasts and may under-represent divergent models
under strict recurrence alone. Because the retinoblastoma contribution to the
final 257-gene pan-cancer feature set depended strongly on isolate rescue rather
than recurrence alone, RBL-specific signals should be interpreted partly as
model-specific structure rather than broadly recurrent signal across multiple
contrasts.

The final 257-gene pan-cancer feature set was not only statistically stable; it
retained interpretable tumour-associated biology. The enrichment analysis linked
disease-specific marker programmes with shared developmental, adhesion, and
migration-related signals in the pan-cancer feature space
(Fig.~\ref{fig:enrichment_top_terms_heatmap}). Curated gene-set labels and TF
motif inferences remain descriptive functional annotations rather than
definitive evidence of direct disease modelling or causal regulatory mechanisms
\citep{Khatri2012PathwayAnalysis}. This conceptual limitation can be compounded
in practice when enrichment workflows use inappropriate background gene
universes or incomplete multiple-testing correction
\citep{Wijesooriya2022Enrichment}.

\subsection{Data comparability and relation to existing transcriptomic mapping frameworks}

Tumour-purity filtering used expression-derived estimates of tumour purity and
immune or stromal admixture \citep{Yoshihara2013}, whereas source harmonisation
used count-aware ComBat-seq modelling for RNA-seq count matrices
\citep{Leek2010BatchEffects,Zhang2020ComBatSeq}.
The purity-filtering and source-harmonisation design reduced admixture- and
source-associated structure in the retained solid-tumour cohorts. Because ESTIMATE infers tumour
purity from expression-derived stromal and immune signatures
\citep{Yoshihara2013}, the purity-filtering design was most appropriate for the
retained solid-tumour cohorts and did not support extending the main ranking and
matching claims to the lymphoma datasets. The weaker
purity--immune/stromal associations in retinoblastoma also indicate that tumour
purity may reflect different biological or technical structure across cohorts,
so purity-filtered similarity results should not be compared as if all cohorts
had identical admixture behaviour. After preprocessing, UMAP projection was
used as a qualitative check of data-origin and lineage structure
\citep{McInnes2020}. The projection showed weaker clustering by data origin and
was not used as a quantitative batch-integration test.

The comparison with existing mapping approaches is methodological rather than a
direct benchmark. Celligner aligns tumour and cell-line transcriptomes using
contrastive PCA and mutual-nearest-neighbour correction
\citep{Warren2021GlobalAlignment}, whereas MFmap uses a semi-supervised
generative latent-space model integrating expression, copy-number, and mutation
features to relate cell lines to tumours and cancer subtypes
\citep{Zhang2021MFmap}. In contrast, the present framework preserves the
curated expression feature space and evaluates bidirectional,
multi-representation similarity rankings and patient-referenced graph
structure. Classification models such as OncoNPC address a different task,
using targeted sequencing-derived mutation features to predict tumour origin
rather than prioritising cell-line models by continuous transcriptomic
similarity \citep{Moon2023CUP}.

\subsection{Limitations and future validation}

The findings should be interpreted within the constraints of the data and
validation design. The current study was restricted to three solid-tumour
lineages and depended on available public bulk RNA-seq cohorts. The findings
are sensitive to threshold and feature--distance choices, and the UMAP
projection served only as a visual aid rather than quantitative integration
evidence. More broadly, transcriptomic similarity does not establish drug
response, essentiality, proteomic state, clonal composition,
microenvironmental interaction, or xenograft behaviour.

The most immediate next steps are testing these rankings against external
tumour cohorts and against pharmacogenomic or dependency resources
\citep{Barretina2012CCLE,Ghandi2019CCLE,Pacini2024DependencyMap}. Proteomic
maps \citep{Goncalves2022ProteomicMap}, single-cell datasets
\citep{Jansky2021NeuroblastomaDevelopment,Yang2021RetinoblastomaSingleCell,Qian2020PanCancerTME},
and patient-derived xenograft models \citep{Liu2023PDXModels} would provide
orthogonal tests of whether transcriptomic proximity aligns with response
phenotypes and \textit{in vivo} model behaviour.

Overall, this study supports transcriptomic model prioritisation as a
patient-referenced ranking problem. The main implication is that cell-line
model selection should be based on transcriptomic proximity to patient tumour
states, with graph structure and marker-gene context used to identify models
whose neighbourhoods require careful interpretation.
